\section{Introduction}
\label{sec:introduction}

Multi-agent systems present unique challenges for reinforcement learning researchers.
Unlike single-agent settings, agents must contend with a non-stationary environment
caused by the simultaneous learning of all other agents \cite{sutton2018reinforcement}.

Despite these challenges, a growing body of work has demonstrated that meaningful
cooperation can emerge without explicit coordination mechanisms. \citet{baker2019emergent}
showed that agents playing hide-and-seek develop sophisticated tool-use behaviours
through purely competitive dynamics.

In this paper, we study three questions:
\begin{enumerate}
  \item Can cooperation emerge from individually rational agents with shared reward?
  \item How does group size affect the speed of emergent coordination?
  \item What structural conditions in the reward function are necessary?
\end{enumerate}

The rest of this paper is organised as follows. Section~\ref{sec:methods} describes
our experimental setup. Section~\ref{sec:conclusion} discusses implications.
